% Options for packages loaded elsewhere
\PassOptionsToPackage{unicode}{hyperref}
\PassOptionsToPackage{hyphens}{url}
\documentclass[
]{book}
\usepackage{xcolor}
\usepackage{amsmath,amssymb}
\setcounter{secnumdepth}{5}
\usepackage{iftex}
\ifPDFTeX
  \usepackage[T1]{fontenc}
  \usepackage[utf8]{inputenc}
  \usepackage{textcomp} % provide euro and other symbols
\else % if luatex or xetex
  \usepackage{unicode-math} % this also loads fontspec
  \defaultfontfeatures{Scale=MatchLowercase}
  \defaultfontfeatures[\rmfamily]{Ligatures=TeX,Scale=1}
\fi
\usepackage{lmodern}
\ifPDFTeX\else
  % xetex/luatex font selection
\fi
% Use upquote if available, for straight quotes in verbatim environments
\IfFileExists{upquote.sty}{\usepackage{upquote}}{}
\IfFileExists{microtype.sty}{% use microtype if available
  \usepackage[]{microtype}
  \UseMicrotypeSet[protrusion]{basicmath} % disable protrusion for tt fonts
}{}
\makeatletter
\@ifundefined{KOMAClassName}{% if non-KOMA class
  \IfFileExists{parskip.sty}{%
    \usepackage{parskip}
  }{% else
    \setlength{\parindent}{0pt}
    \setlength{\parskip}{6pt plus 2pt minus 1pt}}
}{% if KOMA class
  \KOMAoptions{parskip=half}}
\makeatother
\usepackage{color}
\usepackage{fancyvrb}
\newcommand{\VerbBar}{|}
\newcommand{\VERB}{\Verb[commandchars=\\\{\}]}
\DefineVerbatimEnvironment{Highlighting}{Verbatim}{commandchars=\\\{\}}
% Add ',fontsize=\small' for more characters per line
\usepackage{framed}
\definecolor{shadecolor}{RGB}{248,248,248}
\newenvironment{Shaded}{\begin{snugshade}}{\end{snugshade}}
\newcommand{\AlertTok}[1]{\textcolor[rgb]{0.94,0.16,0.16}{#1}}
\newcommand{\AnnotationTok}[1]{\textcolor[rgb]{0.56,0.35,0.01}{\textbf{\textit{#1}}}}
\newcommand{\AttributeTok}[1]{\textcolor[rgb]{0.13,0.29,0.53}{#1}}
\newcommand{\BaseNTok}[1]{\textcolor[rgb]{0.00,0.00,0.81}{#1}}
\newcommand{\BuiltInTok}[1]{#1}
\newcommand{\CharTok}[1]{\textcolor[rgb]{0.31,0.60,0.02}{#1}}
\newcommand{\CommentTok}[1]{\textcolor[rgb]{0.56,0.35,0.01}{\textit{#1}}}
\newcommand{\CommentVarTok}[1]{\textcolor[rgb]{0.56,0.35,0.01}{\textbf{\textit{#1}}}}
\newcommand{\ConstantTok}[1]{\textcolor[rgb]{0.56,0.35,0.01}{#1}}
\newcommand{\ControlFlowTok}[1]{\textcolor[rgb]{0.13,0.29,0.53}{\textbf{#1}}}
\newcommand{\DataTypeTok}[1]{\textcolor[rgb]{0.13,0.29,0.53}{#1}}
\newcommand{\DecValTok}[1]{\textcolor[rgb]{0.00,0.00,0.81}{#1}}
\newcommand{\DocumentationTok}[1]{\textcolor[rgb]{0.56,0.35,0.01}{\textbf{\textit{#1}}}}
\newcommand{\ErrorTok}[1]{\textcolor[rgb]{0.64,0.00,0.00}{\textbf{#1}}}
\newcommand{\ExtensionTok}[1]{#1}
\newcommand{\FloatTok}[1]{\textcolor[rgb]{0.00,0.00,0.81}{#1}}
\newcommand{\FunctionTok}[1]{\textcolor[rgb]{0.13,0.29,0.53}{\textbf{#1}}}
\newcommand{\ImportTok}[1]{#1}
\newcommand{\InformationTok}[1]{\textcolor[rgb]{0.56,0.35,0.01}{\textbf{\textit{#1}}}}
\newcommand{\KeywordTok}[1]{\textcolor[rgb]{0.13,0.29,0.53}{\textbf{#1}}}
\newcommand{\NormalTok}[1]{#1}
\newcommand{\OperatorTok}[1]{\textcolor[rgb]{0.81,0.36,0.00}{\textbf{#1}}}
\newcommand{\OtherTok}[1]{\textcolor[rgb]{0.56,0.35,0.01}{#1}}
\newcommand{\PreprocessorTok}[1]{\textcolor[rgb]{0.56,0.35,0.01}{\textit{#1}}}
\newcommand{\RegionMarkerTok}[1]{#1}
\newcommand{\SpecialCharTok}[1]{\textcolor[rgb]{0.81,0.36,0.00}{\textbf{#1}}}
\newcommand{\SpecialStringTok}[1]{\textcolor[rgb]{0.31,0.60,0.02}{#1}}
\newcommand{\StringTok}[1]{\textcolor[rgb]{0.31,0.60,0.02}{#1}}
\newcommand{\VariableTok}[1]{\textcolor[rgb]{0.00,0.00,0.00}{#1}}
\newcommand{\VerbatimStringTok}[1]{\textcolor[rgb]{0.31,0.60,0.02}{#1}}
\newcommand{\WarningTok}[1]{\textcolor[rgb]{0.56,0.35,0.01}{\textbf{\textit{#1}}}}
\usepackage{longtable,booktabs,array}
\usepackage{calc} % for calculating minipage widths
% Correct order of tables after \paragraph or \subparagraph
\usepackage{etoolbox}
\makeatletter
\patchcmd\longtable{\par}{\if@noskipsec\mbox{}\fi\par}{}{}
\makeatother
% Allow footnotes in longtable head/foot
\IfFileExists{footnotehyper.sty}{\usepackage{footnotehyper}}{\usepackage{footnote}}
\makesavenoteenv{longtable}
\usepackage{graphicx}
\makeatletter
\newsavebox\pandoc@box
\newcommand*\pandocbounded[1]{% scales image to fit in text height/width
  \sbox\pandoc@box{#1}%
  \Gscale@div\@tempa{\textheight}{\dimexpr\ht\pandoc@box+\dp\pandoc@box\relax}%
  \Gscale@div\@tempb{\linewidth}{\wd\pandoc@box}%
  \ifdim\@tempb\p@<\@tempa\p@\let\@tempa\@tempb\fi% select the smaller of both
  \ifdim\@tempa\p@<\p@\scalebox{\@tempa}{\usebox\pandoc@box}%
  \else\usebox{\pandoc@box}%
  \fi%
}
% Set default figure placement to htbp
\def\fps@figure{htbp}
\makeatother
\setlength{\emergencystretch}{3em} % prevent overfull lines
\providecommand{\tightlist}{%
  \setlength{\itemsep}{0pt}\setlength{\parskip}{0pt}}
\usepackage[]{natbib}
\bibliographystyle{apalike}
\usepackage{booktabs}
\usepackage{bookmark}
\IfFileExists{xurl.sty}{\usepackage{xurl}}{} % add URL line breaks if available
\urlstyle{same}
\hypersetup{
  pdftitle={Fundamental Statistics},
  pdfauthor={Taryn Axelsen},
  hidelinks,
  pdfcreator={LaTeX via pandoc}}

\title{Fundamental Statistics}
\author{Taryn Axelsen}
\date{Last Updated: 2026-03-10}

\begin{document}
\maketitle

{
\setcounter{tocdepth}{1}
\tableofcontents
}
\newcommand{\cms}{\,\text{cm}}
\newcommand{\dLs}{\,\text{dL}}
\newcommand{\xdLs}{\text{dL}}

\newcommand{\fmols}{\,\text{fmol}}
\newcommand{\ft}{\,\text{ft}}
\newcommand{\gs}{\,\text{g}}
\newcommand{\hs}{\,\text{h}}
\newcommand{\xhs}{\text{h}}

\newcommand{\has}{\,\text{ha}}
\newcommand{\xhas}{\text{ha}}

\newcommand{\inches}{\,\text{inches}}
\newcommand{\kgs}{\,\text{kg}}
\newcommand{\kms}{\,\text{km}}
\newcommand{\kWhs}{\,\text{kWh}}
\newcommand{\lbs}{\,\text{lb}}
\newcommand{\Ls}{\,\text{L}}
\newcommand{\xLs}{\text{L}}

\frontmatter

\chapter*{Preface}\label{preface}
\addcontentsline{toc}{chapter}{Preface}

Welcome to Fundamental Statistics!

This Workbook is designed to assist you with the \emph{learning through doing} component of this course. The primary learning resource for the course is the lectures and tutorials, however we do have an accompanying textbook De Veaux, Velleman \& Bock,
Stats: Data and Models, fifth edition. We do not cover the whole text book in this
course or in the order it is presented in the book. You are directed to the relevant
readings from the textbook at the end of each chapter within this Workbook. Once
you have watched the lecture recordings, you should consolidate your understanding
by working through the material and tutorials in this Workbook. The textbook
readings will give you further knowledge and help to really solidify the materials
covered each week.

The tutorials align with the modules covered in this course as listed in the Course
Specification. At the end of each tutorial you will notice a Module Checklist and, in
some cases, a Using SPSS Checklist. Use these checklists to make sure that you have
not missed an important concept as you work through the materials. It's a good idea
to tick them off after completing each module and again at revision time.

Statistics is used regularly to help us gain an understanding of the information
which is available to us on a daily basis. The root of this information is data, and
the need to understand this data is what drives statistics. The goal in this course
is to give you tools to enable you to competently use and interpret statistics. These
statistical tools can be viewed in the same way as everyday household tools.

A toolbox needs to have some basics such as a hammer, pliers, wrench and screwdriver.
Similarly, in this course there are some concepts and skills which are essential
for you to master. We call these Threshold Competencies. In order to pass this course,
you will need to demonstrate that you have full understanding of the \emph{Threshold
Competencies}. Beyond this are what we have labeled Expanded Competencies and
these include what you must know in order to achieve higher than a Pass {[}P{]} in the
course. Material that covers these \emph{Expanded Competencies} will be indicated in this
Workbook by an asterisk *.

The emphasis in Fundamental Statistics is in understanding the basic concepts of
statistical reasoning and learning how to use a statistical software package. There is no intention to focus on the particular needs of one or other discipline. That will
come later in other courses. Basic statistical concepts are generic. Once they are
mastered, learning the language of statistics and specific statistical skills needed for
a particular discipline will not be a major step.

We stick fairly closely to the textbook and its approach throughout this course. It
should be noted that for some modules we have included some Supplementary Notes
in this Workbook. We have included this material because the textbook has not
covered this particular content in a way that we would prefer (Binomial Model and
Chi-square Test of Independence).

The broad objectives of this course are stated below. More detailed overall objectives
are given in the Course Specification. Specific objectives for each module are stated
at the start of each module in the lectures.

Best wishes and good studying!

\section*{Overall Course Objectives}\label{overall-course-objectives}
\addcontentsline{toc}{section}{Overall Course Objectives}

On completion of this course students should be able to:

\begin{itemize}
\item
  explore relationships in data and distinguish between different methods of
  data collection and analysis;
\item
  evaluate and apply a variety of statistical inferential methods to real life situations;
\item
  use a statistical computer package to enter, summarise and analyse data; and
\item
  interpret and communicate the results of statistical analyses for a diverse audience.
\end{itemize}

\section*{How this book was made}\label{how-this-book-was-made}
\addcontentsline{toc}{section}{How this book was made}

This book was made using \textbf{R} \citep{R-base}, and the \textbf{bookdown} package \citep{R-bookdown}, which is based on \href{https://en.wikipedia.org/wiki/Markdown}{Markdown} syntax, using \textbf{knitr} \citep{R-knitr}.

\begin{flushright}
Taryn Axelsen\\
Springfield, Australia
Rachel King
Toowoomba, Australia
Bethany Rognoni
Toowoomba, Australia
\end{flushright}

\mainmatter

\chapter*{About}\label{about}
\addcontentsline{toc}{chapter}{About}

This is a \emph{sample} book written in \textbf{Markdown}. You can use anything that Pandoc's Markdown supports; for example, a math equation \(a^2 + b^2 = c^2\).

\section*{Usage}\label{usage}
\addcontentsline{toc}{section}{Usage}

Each \textbf{bookdown} chapter is an .Rmd file, and each .Rmd file can contain one (and only one) chapter. A chapter \emph{must} start with a first-level heading: \texttt{\#\ A\ good\ chapter}, and can contain one (and only one) first-level heading.

Use second-level and higher headings within chapters like: \texttt{\#\#\ A\ short\ section} or \texttt{\#\#\#\ An\ even\ shorter\ section}.

The \texttt{index.Rmd} file is required, and is also your first book chapter. It will be the homepage when you render the book.

\section*{Render book}\label{render-book}
\addcontentsline{toc}{section}{Render book}

You can render the HTML version of this example book without changing anything:

\begin{enumerate}
\def\labelenumi{\arabic{enumi}.}
\item
  Find the \textbf{Build} pane in the RStudio IDE, and
\item
  Click on \textbf{Build Book}, then select your output format, or select ``All formats'' if you'd like to use multiple formats from the same book source files.
\end{enumerate}

Or build the book from the R console:

\begin{Shaded}
\begin{Highlighting}[]
\NormalTok{bookdown}\SpecialCharTok{::}\FunctionTok{render\_book}\NormalTok{()}
\end{Highlighting}
\end{Shaded}

To render this example to PDF as a \texttt{bookdown::pdf\_book}, you'll need to install XeLaTeX. You are recommended to install TinyTeX (which includes XeLaTeX): \url{https://yihui.org/tinytex/}.

\section*{Preview book}\label{preview-book}
\addcontentsline{toc}{section}{Preview book}

As you work, you may start a local server to live preview this HTML book. This preview will update as you edit the book when you save individual .Rmd files. You can start the server in a work session by using the RStudio add-in ``Preview book'', or from the R console:

\begin{Shaded}
\begin{Highlighting}[]
\NormalTok{bookdown}\SpecialCharTok{::}\FunctionTok{serve\_book}\NormalTok{()}
\end{Highlighting}
\end{Shaded}

\chapter{Summary of Module 1}\label{summary-of-module-1}

Put a summary here.

\section{Review Questions Module 1}\label{review-questions-module-1}

Once you have watched the lecture recordings you should be able to answer the
following \emph{Threshold} questions:

\begin{enumerate}
\def\labelenumi{(\roman{enumi})}
\item
  What is a variable? What is a value of a variable? Give examples of each.
\item
  How can we tell if a variable is quantitative, categorical, ordinal or an identifier?
\item
  What is the difference between a statistic and a parameter? What do they measure
  and how might they be obtained?
\item
  What is the difference between a representative sample and a convenience
  sample? How are they obtained and what are their relative advantages/disadvantages?
\item
  What is a contingency or two-way table? When is it used?
\item
  What is meant by a marginal distribution, a joint distribution and a conditional
  distribution?
\end{enumerate}

\section{Tutorial Questions Module 1}\label{tutorial-questions-module-1}

Note: * Indicates \textbf{Expanded} Competencies

\subsubsection{Question 1: Generating an appropriate sample}\label{question-1-generating-an-appropriate-sample}

\section{Interactive: Sampling (SRS vs Stratified)}\label{interactive-sampling-srs-vs-stratified}

Click the button below to launch the interactive sampling tool.

\subsection{Textbook \#\#\# Readings}\label{textbook-readings}

This module covers material in Chapters 1, 2, 10 and 11 of De Veaux, Velleman \&
Bock, (4th edition)(Chapters 1, 2.1, 3 and 10 of De Veaux, Velleman \& Bock, (5th
edition)).
\#\#\# Reading
For extra clarification on materials in this module, see: De Veaux, Velle-
man \& Bock, (4th edition), Chapters 1, 2, 10 and 11 (De Veaux, Velle-
man \& Bock, (5th edition), Chapters 1, 2.1, 3 and 10).

\bibliography{book.bib,packages.bib}

\end{document}
